\documentclass[11pt,a4paper]{article}
\usepackage[T1]{fontenc}\usepackage[utf8]{inputenc}\usepackage{lmodern}
\usepackage[margin=2.4cm]{geometry}
\usepackage{xcolor}\usepackage{amsmath,amssymb,amsthm,booktabs,longtable,hyperref,enumitem}
\hypersetup{colorlinks=true,linkcolor=blue!50!black,urlcolor=blue!50!black,citecolor=blue!50!black}
\newtheorem{theorem}{Theorem}\newtheorem{lemma}[theorem]{Lemma}\newtheorem{corollary}[theorem]{Corollary}
\theoremstyle{definition}\newtheorem{definition}[theorem]{Definition}
\newtheorem{observation}[theorem]{Computational observation}\newtheorem{question}[theorem]{Open question}
\theoremstyle{remark}\newtheorem{remark}[theorem]{Remark}
\newcommand{\Rep}{\mathrm{Rep}}
\title{Two new counterexamples to a conjecture of Zhi-Wei Sun\\ on $p+2^x+7\cdot 2^y$ (OEIS A157225)}
\author{Jan Jarosz}
\date{26 November 2026}
\begin{document}\maketitle
\begin{abstract}
Sun conjectured in 2009 that every odd integer $m>20$ other than $25$, $31$ and $1983$ can be written as
$m=p+2^x+7\cdot 2^y$ with $p$ a prime congruent to $5$ modulo $6$ and $x,y\ge 1$. The conjecture was disproved in
August 2026 by a single counterexample, $m=1\,433\,987\,797$. We give two more:
\[ m=45\,483\,661\,477\qquad\text{and}\qquad m=85\,305\,551\,767 . \]
Both are proved in Lean~4 with Mathlib, by a certificate that the kernel checks in a few seconds. An exhaustive
sieve finds no other counterexample below $10^{11}$; this last statement is a computation, not a formal proof,
and is labelled as such throughout.
\end{abstract}

\section{Statement}
\begin{definition}
For a positive integer $m$ write $\Rep(m)$ if there are a prime $p\equiv 5\pmod 6$ and integers $x,y\ge 1$ with
$m=p+2^x+7\cdot 2^y$. OEIS A157225 is the sequence $a(n)=\#\{(p,x,y)\}$ for $m=2n-1$, so $a(n)=0$ exactly when
$\Rep(2n-1)$ fails.
\end{definition}

\noindent\textbf{Conjecture (Z.-W. Sun, 24 February 2009 \cite{oeis,sun}).} $a(n)=0$ if and only if $n<11$ or
$n\in\{13,16,992\}$. Equivalently, $\Rep(m)$ holds for every odd $m>20$ except $m=25,31,1983$.
Sun verified this for odd $m<5\cdot 10^{7}$ and Q.-H. Hou for $m<1.5\cdot 10^{8}$ \cite{oeis}.

\medskip\noindent\textbf{Known disproof.} T.~Adamczewski \cite{epoch} reported $a(716\,993\,899)=0$, that is,
$\Rep(1\,433\,987\,797)$ fails. As of 26 November 2026 this is the only counterexample listed on the OEIS page
\cite{oeis}.

\section{Results}
\begin{theorem}[proved; Lean kernel]\label{thm:main}
$\Rep(m)$ fails for
\begin{center}\begin{tabular}{rrl}\toprule
$m$ & $n=(m+1)/2$ & status\\\midrule
$1\,433\,987\,797$ & $716\,993\,899$ & known \cite{epoch}, re-proved here\\
$45\,483\,661\,477$ & $22\,741\,830\,739$ & new\\
$85\,305\,551\,767$ & $42\,652\,775\,884$ & new\\\bottomrule\end{tabular}\end{center}
Hence $a(22\,741\,830\,739)=a(42\,652\,775\,884)=0$.
\end{theorem}
\begin{corollary}[proved]\label{cor}Sun's conjecture for A157225 is false. (This was already known from \cite{epoch}; each of the
three numbers gives an independent proof.)\end{corollary}

\begin{observation}[sieve; not formally verified]\label{obs}
Among odd $m$ with $1983<m<10^{11}$, the three numbers of Theorem~\ref{thm:main} are the only ones for which
$\Rep(m)$ fails. In particular $1\,433\,987\,797$ is the least counterexample and the two new numbers are the
second and the third.
\end{observation}

\noindent\textbf{What is and is not claimed.}
\begin{itemize}[leftmargin=1.4em,itemsep=1pt]
\item \emph{Counterexamples (theorems):} the three numbers in Theorem~\ref{thm:main}. Nothing here is a mere candidate.
\item \emph{Computed, not proved:} completeness below $10^{11}$ (Observation~\ref{obs}). It rests on our sieve
program and on the hardware. It was run twice with identical output and it reproduces Sun's list of small
exceptions and Adamczewski's number.
\item \emph{Novelty:} ``new'' means not on the OEIS page \cite{oeis} on 26 November 2026 and unknown to us from
any other source. We have not made an exhaustive literature search.
\item \emph{Not claimed:} that there are infinitely many counterexamples (Question~\ref{q}), or anything about the
sister conjectures A157237 and A157218, which we checked up to $10^{11}$ without finding a counterexample.
\end{itemize}

\section{Proof of Theorem~\ref{thm:main}}
Fix $m$. Call a pair $(x,y)$ with $x,y\ge1$ and $2^x+7\cdot2^y<m$ a \emph{cell}, and put $r=m-2^x-7\cdot 2^y$.
$\Rep(m)$ holds iff for some cell $r$ is a prime $\equiv5\pmod 6$. There are finitely many cells
($x,y<\log_2 m$), so it suffices to exhibit for each cell either $r\not\equiv 5\pmod 6$ or a divisor $d$ with $1<d<r$,
$d\mid r$.

\begin{lemma}[parity filter]\label{lem}
Let $m\equiv 1\pmod 6$. If $r\equiv 5\pmod 6$ then $x$ and $y$ are both even.
\end{lemma}
\begin{proof}
$2^x\equiv 2$ or $4\pmod 6$ according as $x$ is odd or even, and $7\cdot2^y\equiv 2^y$. We need
$2^x+7\cdot 2^y\equiv m-5\equiv 2\pmod 6$, and among $2+2$, $2+4$, $4+4$ only $4+4\equiv 2$.
\end{proof}
All three numbers are $\equiv1\pmod 6$, so about three quarters of the cells die at once. The remaining cells
are killed by explicit divisors:
\begin{center}\begin{tabular}{rrrrrr}\toprule
$m$ & cells & $r\not\equiv5\ (6)$ & composite $r$ & distinct divisors used & largest\\\midrule
$1\,433\,987\,797$ & 807 & 613 & 194 & 40 & 13\,313\\
$45\,483\,661\,477$ & 1117 & 846 & 271 & 60 & 163\,883\\
$85\,305\,551\,767$ & 1185 & 898 & 287 & 66 & 292\,057\\\bottomrule\end{tabular}\end{center}
The divisor of every surviving cell is listed in the Appendix, so the theorem can be checked by
$271+287$ divisions (plus $194$ for the known number). \qed

\paragraph{Formal proof.} The file \texttt{SunA157225.lean} (Lean 4.35.0, Mathlib) defines
$\Rep$ with Mathlib's \texttt{Nat.Prime}, a Boolean function \texttt{check m N L} that tests all cells with $x,y<N$
against a divisor list $L$, and proves the soundness lemma
\[\texttt{not\_rep\_of\_check}:\quad m<2^N\ \wedge\ \texttt{check m N L}=\texttt{true}\ \Longrightarrow\ \neg\,\Rep(m).\]
The three theorems \texttt{not\_rep\_1433987797}, \texttt{not\_rep\_45483661477}, \texttt{not\_rep\_85305551767}
apply it with $N=31,36,37$; the evaluation of \texttt{check} is done by the Lean kernel (\texttt{decide +kernel}).
There is no \texttt{sorry} and no \texttt{native\_decide}; the axioms used are only \texttt{propext} and
\texttt{Quot.sound}. Corollary~\ref{cor} is \texttt{sunConjecture\_false}. The whole file compiles in about 5 seconds.

Observation~\ref{obs} appears in the file only as a named hypothesis \texttt{SieveComplete : Prop}. Two
conditional theorems use it as an explicit assumption: \texttt{least\_of\_sieve} (then $1\,433\,987\,797$ is the least
counterexample) and \texttt{counter\_below\_of\_sieve} (then the set of counterexamples below $10^{11}$ is exactly
our three numbers).

\section{Method}
\paragraph{Search.} A bit-sieve. For each residue $R\in\{1,5\}$ a bitmap stores which $6k+R<B$ are prime.
For a fixed class $m\equiv c\pmod 6$, each term $t=2^x+7\cdot 2^y$ with $t\equiv c-5 \pmod 6$ turns the question
``is $m-t$ prime?'' into a shift of the bitmap by $(t-c+5)/6$ positions. OR-ing the shifted bitmaps over all $t$
gives, 64 values of $m$ per machine word, the set of representable $m$; zero bits are the exceptions. With
$B=10^{11}$ this takes about 75 seconds on 4 cores and 2\,GB of memory (\texttt{sieve2k.c}). The complete output
for odd $m<10^{11}$ is
\[1,3,5,\dots,19,\ 25,\ 31,\ 1983,\ 1\,433\,987\,797,\ 45\,483\,661\,477,\ 85\,305\,551\,767 .\]
\paragraph{Independent check.} Each of the three numbers was re-examined by a separate Python script (SymPy
primality and factoring, no shared code with the sieve): every cell was enumerated and its $r$ shown to be
$\not\equiv 5\pmod 6$ or composite. The same script recovers exactly Sun's exceptions for $m<2\cdot10^{5}$.
\paragraph{How the candidates were ranked.} The search was guided by a ``clock'' picture: a small prime $q$ with
small multiplicative order of $2$ kills a fixed periodic pattern of cells once $m \bmod q$ is fixed. The clock $3$
gives Lemma~\ref{lem}; for our numbers the next clocks $5$ and $7$ kill $72$ and about $71$ of the roughly $280$
surviving cells each, and $13$ and $17$ another $15$ to $26$. This explains why the counterexamples exist but it is
a heuristic: we found no full covering system, so it proves nothing about
further counterexamples.

\begin{question}\label{q}
Are there infinitely many odd $m$ for which $\Rep(m)$ fails? The number of cells surviving the parity filter grows
like $(\log_2 m)^2/4$ while the chance that a given $r$ is prime decays only like $1/\log m$, so
counterexamples should become rarer; without a covering congruence argument (as in Crocker \cite{crocker} for
$p+2^x+2^y$) we see no way to decide this. The next counterexample, if any, exceeds $10^{11}$.
\end{question}

\section{Files}
\begin{itemize}[leftmargin=1.4em,itemsep=1pt]
\item \texttt{SunA157225.lean}: definitions, checker, the three theorems, hypotheses and conditional theorems.
\item \texttt{certificate\_}$m$\texttt{.csv} for the three $m$: one line $(x,y,r,d)$ per surviving cell, $d$ the least prime factor of $r$.
\item \texttt{verify.py}: the independent check. \texttt{sieve2k.c}: the sieve
(\texttt{./sieve2k 5 100000000000 7}).
\end{itemize}

\begin{thebibliography}{9}
\bibitem{oeis} OEIS Foundation, sequence A157225, \url{https://oeis.org/A157225} (consulted 26 November 2026).
\bibitem{sun} Z.-W. Sun, \emph{Mixed sums of primes and other terms}, arXiv:0901.3075 (2009).
\bibitem{epoch} T. Adamczewski, \emph{OEIS Open: How many conjectures can language models turn into theorems?}, arXiv:2608.11941 (2026).
\bibitem{crocker} R. Crocker, \emph{On a sum of a prime and two powers of two}, Pacific J. Math. 36 (1971), 103--107.
\end{thebibliography}

\appendix
\section{Certificates}
Each entry is $(x,y)\ d$: for the cell $(x,y)$ with $x,y$ both even, $d$ is the least prime factor of
$r=m-2^x-7\cdot2^y$ (and $d<r$). Cells not listed have an odd $x$ or $y$ and are excluded by Lemma~\ref{lem}.
\small
\subsection*{$m=45\,483\,661\,477$ \quad (271 cells; $x\le 34$, $y\le 32$)}
Divisors used: 5, 7, 11, 13, 17, 19, 23, 29, 31, 37, 47, 53, 61, 67, 79, 83, 89, 97, 101, 107, 113, 131, 139, 163, 173, 179, 223, 227, 239, 241, 269, 277, 281, 313, 461, 479, 509, 563, 1163, 1249, 1433, 1447, 1733, 1747, 2113, 2389, 4159, 5653, 6619, 13163, 20897, 29131, 33563, 45869, 51229, 58997, 81181, 81517, 82571, 163883.
\begin{longtable}{@{}llllll@{}}(2,2)\,5 & (2,4)\,223 & (2,6)\,5 & (2,8)\,17 & (2,10)\,5 & (2,12)\,37 \\
(2,14)\,5 & (2,16)\,11 & (2,18)\,5 & (2,20)\,23 & (2,22)\,5 & (2,24)\,17 \\
(2,26)\,5 & (2,28)\,33563 & (2,30)\,5 & (2,32)\,17 & (4,2)\,7 & (4,4)\,7 \\
(4,6)\,7 & (4,8)\,7 & (4,10)\,7 & (4,12)\,7 & (4,14)\,7 & (4,16)\,7 \\
(4,18)\,7 & (4,20)\,7 & (4,22)\,7 & (4,24)\,7 & (4,26)\,7 & (4,28)\,7 \\
(4,30)\,7 & (4,32)\,7 & (6,2)\,5 & (6,4)\,13 & (6,6)\,5 & (6,8)\,53 \\
(6,10)\,5 & (6,12)\,29 & (6,14)\,5 & (6,16)\,13 & (6,18)\,5 & (6,20)\,19 \\
(6,22)\,5 & (6,24)\,23 & (6,26)\,5 & (6,28)\,13 & (6,30)\,5 & (6,32)\,113 \\
(8,2)\,1733 & (8,4)\,17 & (8,6)\,37 & (8,8)\,13 & (8,10)\,61 & (8,12)\,17 \\
(8,14)\,277 & (8,16)\,4159 & (8,18)\,19 & (8,20)\,13 & (8,22)\,269 & (8,24)\,97 \\
(8,26)\,23 & (8,28)\,17 & (8,30)\,45869 & (8,32)\,13 & (10,2)\,5 & (10,4)\,7 \\
(10,6)\,5 & (10,8)\,7 & (10,10)\,5 & (10,12)\,7 & (10,14)\,5 & (10,16)\,7 \\
(10,18)\,5 & (10,20)\,7 & (10,22)\,5 & (10,24)\,7 & (10,26)\,5 & (10,28)\,7 \\
(10,30)\,5 & (10,32)\,7 & (12,2)\,82571 & (12,4)\,281 & (12,6)\,11 & (12,8)\,23 \\
(12,10)\,19 & (12,12)\,461 & (12,14)\,2389 & (12,16)\,11 & (12,18)\,13 & (12,20)\,31 \\
(12,22)\,81181 & (12,24)\,131 & (12,26)\,11 & (12,28)\,19 & (12,30)\,13 & (12,32)\,37 \\
(14,2)\,5 & (14,4)\,1433 & (14,6)\,5 & (14,8)\,6619 & (14,10)\,5 & (14,12)\,51229 \\
(14,14)\,5 & (14,16)\,239 & (14,18)\,5 & (14,20)\,11 & (14,22)\,5 & (14,24)\,113 \\
(14,26)\,5 & (14,28)\,1447 & (14,30)\,5 & (14,32)\,23 & (16,2)\,7 & (16,4)\,7 \\
(16,6)\,7 & (16,8)\,7 & (16,10)\,7 & (16,12)\,7 & (16,14)\,7 & (16,16)\,7 \\
(16,18)\,7 & (16,20)\,7 & (16,22)\,7 & (16,24)\,7 & (16,26)\,7 & (16,28)\,7 \\
(16,30)\,7 & (16,32)\,7 & (18,2)\,5 & (18,4)\,13 & (18,6)\,5 & (18,8)\,17 \\
(18,10)\,5 & (18,12)\,173 & (18,14)\,5 & (18,16)\,13 & (18,18)\,5 & (18,20)\,227 \\
(18,22)\,5 & (18,24)\,17 & (18,26)\,5 & (18,28)\,13 & (18,30)\,5 & (18,32)\,17 \\
(20,2)\,163883 & (20,4)\,67 & (20,6)\,179 & (20,8)\,13 & (20,10)\,163 & (20,12)\,1163 \\
(20,14)\,58997 & (20,16)\,23 & (20,18)\,107 & (20,20)\,13 & (20,22)\,20897 & (20,24)\,313 \\
(20,26)\,1747 & (20,28)\,131 & (20,30)\,241 & (20,32)\,13 & (22,2)\,5 & (22,4)\,7 \\
(22,6)\,5 & (22,8)\,7 & (22,10)\,5 & (22,12)\,7 & (22,14)\,5 & (22,16)\,7 \\
(22,18)\,5 & (22,20)\,7 & (22,22)\,5 & (22,24)\,7 & (22,26)\,5 & (22,28)\,7 \\
(22,30)\,5 & (22,32)\,7 & (24,2)\,19 & (24,4)\,17 & (24,6)\,13 & (24,8)\,61 \\
(24,10)\,11 & (24,12)\,17 & (24,14)\,479 & (24,16)\,509 & (24,18)\,13 & (24,20)\,11 \\
(24,22)\,29 & (24,24)\,139 & (24,26)\,47 & (24,28)\,17 & (24,30)\,11 & (24,32)\,29131 \\
(26,2)\,5 & (26,4)\,1733 & (26,6)\,5 & (26,8)\,17 & (26,10)\,5 & (26,12)\,1249 \\
(26,14)\,5 & (26,16)\,17 & (26,18)\,5 & (26,20)\,61 & (26,22)\,5 & (26,24)\,17 \\
(26,26)\,5 & (26,28)\,13163 & (26,30)\,5 & (26,32)\,17 & (28,2)\,7 & (28,4)\,7 \\
(28,6)\,7 & (28,8)\,7 & (28,10)\,7 & (28,12)\,7 & (28,14)\,7 & (28,16)\,7 \\
(28,18)\,7 & (28,20)\,7 & (28,22)\,7 & (28,24)\,7 & (28,26)\,7 & (28,28)\,7 \\
(28,30)\,7 & (28,32)\,7 & (30,2)\,5 & (30,4)\,13 & (30,6)\,5 & (30,8)\,79 \\
(30,10)\,5 & (30,12)\,81517 & (30,14)\,5 & (30,16)\,13 & (30,18)\,5 & (30,20)\,89 \\
(30,22)\,5 & (30,24)\,83 & (30,26)\,5 & (30,28)\,13 & (30,30)\,5 & (30,32)\,131 \\
(32,2)\,101 & (32,4)\,17 & (32,6)\,11 & (32,8)\,13 & (32,10)\,31 & (32,12)\,17 \\
(32,14)\,563 & (32,16)\,11 & (32,18)\,5653 & (32,20)\,13 & (32,22)\,83 & (32,24)\,2113 \\
(32,26)\,11 & (32,28)\,17 & (32,30)\,31 & (32,32)\,13 & (34,2)\,5 & (34,4)\,7 \\
(34,6)\,5 & (34,8)\,7 & (34,10)\,5 & (34,12)\,7 & (34,14)\,5 & (34,16)\,7 \\
(34,18)\,5 & (34,20)\,7 & (34,22)\,5 & (34,24)\,7 & (34,26)\,5 & (34,28)\,7 \\
(34,30)\,5 \\
\end{longtable}
\subsection*{$m=85\,305\,551\,767$ \quad (287 cells; $x\le 36$, $y\le 32$)}
Divisors used: 5, 7, 11, 13, 17, 19, 23, 29, 31, 37, 41, 43, 47, 53, 67, 79, 83, 89, 101, 103, 107, 109, 127, 131, 151, 157, 173, 179, 181, 211, 229, 241, 271, 277, 467, 757, 797, 821, 1867, 2011, 2351, 2797, 3533, 3593, 4229, 4483, 4673, 4931, 5261, 5849, 7901, 8167, 8839, 9437, 11251, 13229, 14753, 20483, 21961, 35053, 68059, 85469, 151531, 163433, 231131, 292057.
\begin{longtable}{@{}llllll@{}}(2,2)\,5 & (2,4)\,41 & (2,6)\,5 & (2,8)\,21961 & (2,10)\,5 & (2,12)\,467 \\
(2,14)\,5 & (2,16)\,271 & (2,18)\,5 & (2,20)\,173 & (2,22)\,5 & (2,24)\,41 \\
(2,26)\,5 & (2,28)\,19 & (2,30)\,5 & (2,32)\,4673 & (4,2)\,7 & (4,4)\,7 \\
(4,6)\,7 & (4,8)\,7 & (4,10)\,7 & (4,12)\,7 & (4,14)\,7 & (4,16)\,7 \\
(4,18)\,7 & (4,20)\,7 & (4,22)\,7 & (4,24)\,7 & (4,26)\,7 & (4,28)\,7 \\
(4,30)\,7 & (4,32)\,7 & (6,2)\,5 & (6,4)\,13 & (6,6)\,5 & (6,8)\,37 \\
(6,10)\,5 & (6,12)\,17 & (6,14)\,5 & (6,16)\,13 & (6,18)\,5 & (6,20)\,17 \\
(6,22)\,5 & (6,24)\,2351 & (6,26)\,5 & (6,28)\,13 & (6,30)\,5 & (6,32)\,43 \\
(8,2)\,23 & (8,4)\,8839 & (8,6)\,157 & (8,8)\,13 & (8,10)\,31 & (8,12)\,107 \\
(8,14)\,29 & (8,16)\,4483 & (8,18)\,292057 & (8,20)\,13 & (8,22)\,131 & (8,24)\,23 \\
(8,26)\,9437 & (8,28)\,231131 & (8,30)\,31 & (8,32)\,13 & (10,2)\,5 & (10,4)\,7 \\
(10,6)\,5 & (10,8)\,7 & (10,10)\,5 & (10,12)\,7 & (10,14)\,5 & (10,16)\,7 \\
(10,18)\,5 & (10,20)\,7 & (10,22)\,5 & (10,24)\,7 & (10,26)\,5 & (10,28)\,7 \\
(10,30)\,5 & (10,32)\,7 & (12,2)\,67 & (12,4)\,19 & (12,6)\,13 & (12,8)\,17 \\
(12,10)\,7901 & (12,12)\,101 & (12,14)\,83 & (12,16)\,17 & (12,18)\,13 & (12,20)\,11251 \\
(12,22)\,19 & (12,24)\,17 & (12,26)\,41 & (12,28)\,163433 & (12,30)\,13 & (12,32)\,17 \\
(14,2)\,5 & (14,4)\,11 & (14,6)\,5 & (14,8)\,47 & (14,10)\,5 & (14,12)\,17 \\
(14,14)\,5 & (14,16)\,13229 & (14,18)\,5 & (14,20)\,17 & (14,22)\,5 & (14,24)\,11 \\
(14,26)\,5 & (14,28)\,17 & (14,30)\,5 & (14,32)\,2011 & (16,2)\,7 & (16,4)\,7 \\
(16,6)\,7 & (16,8)\,7 & (16,10)\,7 & (16,12)\,7 & (16,14)\,7 & (16,16)\,7 \\
(16,18)\,7 & (16,20)\,7 & (16,22)\,7 & (16,24)\,7 & (16,26)\,7 & (16,28)\,7 \\
(16,30)\,7 & (16,32)\,7 & (18,2)\,5 & (18,4)\,13 & (18,6)\,5 & (18,8)\,151 \\
(18,10)\,5 & (18,12)\,103 & (18,14)\,5 & (18,16)\,13 & (18,18)\,5 & (18,20)\,31 \\
(18,22)\,5 & (18,24)\,151531 & (18,26)\,5 & (18,28)\,13 & (18,30)\,5 & (18,32)\,127 \\
(20,2)\,11 & (20,4)\,43 & (20,6)\,31 & (20,8)\,13 & (20,10)\,19 & (20,12)\,11 \\
(20,14)\,8167 & (20,16)\,17 & (20,18)\,43 & (20,20)\,13 & (20,22)\,11 & (20,24)\,17 \\
(20,26)\,31 & (20,28)\,19 & (20,30)\,1867 & (20,32)\,11 & (22,2)\,5 & (22,4)\,7 \\
(22,6)\,5 & (22,8)\,7 & (22,10)\,5 & (22,12)\,7 & (22,14)\,5 & (22,16)\,7 \\
(22,18)\,5 & (22,20)\,7 & (22,22)\,5 & (22,24)\,7 & (22,26)\,5 & (22,28)\,7 \\
(22,30)\,5 & (22,32)\,7 & (24,2)\,20483 & (24,4)\,11 & (24,6)\,13 & (24,8)\,2797 \\
(24,10)\,131 & (24,12)\,85469 & (24,14)\,11 & (24,16)\,47 & (24,18)\,13 & (24,20)\,229 \\
(24,22)\,3593 & (24,24)\,11 & (24,26)\,101 & (24,28)\,79 & (24,30)\,13 & (24,32)\,35053 \\
(26,2)\,5 & (26,4)\,103 & (26,6)\,5 & (26,8)\,79 & (26,10)\,5 & (26,12)\,757 \\
(26,14)\,5 & (26,16)\,89 & (26,18)\,5 & (26,20)\,67 & (26,22)\,5 & (26,24)\,29 \\
(26,26)\,5 & (26,28)\,277 & (26,30)\,5 & (26,32)\,23 & (28,2)\,7 & (28,4)\,7 \\
(28,6)\,7 & (28,8)\,7 & (28,10)\,7 & (28,12)\,7 & (28,14)\,7 & (28,16)\,7 \\
(28,18)\,7 & (28,20)\,7 & (28,22)\,7 & (28,24)\,7 & (28,26)\,7 & (28,28)\,7 \\
(28,30)\,7 & (28,32)\,7 & (30,2)\,5 & (30,4)\,13 & (30,6)\,5 & (30,8)\,797 \\
(30,10)\,5 & (30,12)\,11 & (30,14)\,5 & (30,16)\,13 & (30,18)\,5 & (30,20)\,17 \\
(30,22)\,5 & (30,24)\,23 & (30,26)\,5 & (30,28)\,13 & (30,30)\,5 & (30,32)\,11 \\
(32,2)\,4229 & (32,4)\,127 & (32,6)\,41 & (32,8)\,13 & (32,10)\,5261 & (32,12)\,4931 \\
(32,14)\,821 & (32,16)\,29 & (32,18)\,53 & (32,20)\,13 & (32,22)\,37 & (32,24)\,241 \\
(32,26)\,41 & (32,28)\,5849 & (32,30)\,181 & (32,32)\,13 & (34,2)\,5 & (34,4)\,7 \\
(34,6)\,5 & (34,8)\,7 & (34,10)\,5 & (34,12)\,7 & (34,14)\,5 & (34,16)\,7 \\
(34,18)\,5 & (34,20)\,7 & (34,22)\,5 & (34,24)\,7 & (34,26)\,5 & (34,28)\,7 \\
(34,30)\,5 & (34,32)\,7 & (36,2)\,14753 & (36,4)\,67 & (36,6)\,13 & (36,8)\,17 \\
(36,10)\,179 & (36,12)\,68059 & (36,14)\,29 & (36,16)\,17 & (36,18)\,13 & (36,20)\,131 \\
(36,22)\,109 & (36,24)\,17 & (36,26)\,3533 & (36,28)\,211 & (36,30)\,13 \\
\end{longtable}
\subsection*{$m=1\,433\,987\,797$ \quad (194 cells; known counterexample)}
Divisors used: 5, 7, 11, 13, 17, 19, 23, 31, 37, 41, 43, 47, 53, 59, 79, 83, 89, 103, 109, 113, 157, 163, 239, 281, 359, 367, 373, 439, 479, 499, 631, 659, 877, 1753, 2543, 2617, 3167, 4219, 10159, 13313.
\begin{longtable}{@{}llllll@{}}(2,2)\,5 & (2,4)\,7 & (2,6)\,5 & (2,8)\,7 & (2,10)\,5 & (2,12)\,7 \\
(2,14)\,5 & (2,16)\,7 & (2,18)\,5 & (2,20)\,7 & (2,22)\,5 & (2,24)\,7 \\
(2,26)\,5 & (4,2)\,11 & (4,4)\,41 & (4,6)\,239 & (4,8)\,17 & (4,10)\,13 \\
(4,12)\,11 & (4,14)\,2543 & (4,16)\,17 & (4,18)\,439 & (4,20)\,23 & (4,22)\,11 \\
(4,24)\,17 & (4,26)\,89 & (6,2)\,5 & (6,4)\,11 & (6,6)\,5 & (6,8)\,31 \\
(6,10)\,5 & (6,12)\,17 & (6,14)\,5 & (6,16)\,23 & (6,18)\,5 & (6,20)\,17 \\
(6,22)\,5 & (6,24)\,11 & (6,26)\,5 & (8,2)\,7 & (8,4)\,7 & (8,6)\,7 \\
(8,8)\,7 & (8,10)\,7 & (8,12)\,7 & (8,14)\,7 & (8,16)\,7 & (8,18)\,7 \\
(8,20)\,7 & (8,22)\,7 & (8,24)\,7 & (8,26)\,7 & (10,2)\,5 & (10,4)\,43 \\
(10,6)\,5 & (10,8)\,11 & (10,10)\,5 & (10,12)\,19 & (10,14)\,5 & (10,16)\,1753 \\
(10,18)\,5 & (10,20)\,13 & (10,22)\,5 & (10,24)\,877 & (10,26)\,5 & (12,2)\,113 \\
(12,4)\,79 & (12,6)\,4219 & (12,8)\,17 & (12,10)\,11 & (12,12)\,13 & (12,14)\,23 \\
(12,16)\,17 & (12,18)\,19 & (12,20)\,11 & (12,22)\,659 & (12,24)\,13 & (12,26)\,3167 \\
(14,2)\,5 & (14,4)\,7 & (14,6)\,5 & (14,8)\,7 & (14,10)\,5 & (14,12)\,7 \\
(14,14)\,5 & (14,16)\,7 & (14,18)\,5 & (14,20)\,7 & (14,22)\,5 & (14,24)\,7 \\
(14,26)\,5 & (16,2)\,41 & (16,4)\,11 & (16,6)\,103 & (16,8)\,31 & (16,10)\,13 \\
(16,12)\,367 & (16,14)\,11 & (16,16)\,37 & (16,18)\,31 & (16,20)\,479 & (16,22)\,13 \\
(16,24)\,11 & (16,26)\,10159 & (18,2)\,5 & (18,4)\,19 & (18,6)\,5 & (18,8)\,47 \\
(18,10)\,5 & (18,12)\,163 & (18,14)\,5 & (18,16)\,11 & (18,18)\,5 & (18,20)\,359 \\
(18,22)\,5 & (18,24)\,31 & (18,26)\,5 & (20,2)\,7 & (20,4)\,7 & (20,6)\,7 \\
(20,8)\,7 & (20,10)\,7 & (20,12)\,7 & (20,14)\,7 & (20,16)\,7 & (20,18)\,7 \\
(20,20)\,7 & (20,22)\,7 & (20,24)\,7 & (20,26)\,7 & (22,2)\,5 & (22,4)\,17 \\
(22,6)\,5 & (22,8)\,13 & (22,10)\,5 & (22,12)\,17 & (22,14)\,5 & (22,16)\,2617 \\
(22,18)\,5 & (22,20)\,11 & (22,22)\,5 & (22,24)\,53 & (22,26)\,5 & (24,2)\,11 \\
(24,4)\,41 & (24,6)\,19 & (24,8)\,13313 & (24,10)\,499 & (24,12)\,11 & (24,14)\,37 \\
(24,16)\,157 & (24,18)\,43 & (24,20)\,59 & (24,22)\,11 & (24,24)\,13 & (24,26)\,877 \\
(26,2)\,5 & (26,4)\,7 & (26,6)\,5 & (26,8)\,7 & (26,10)\,5 & (26,12)\,7 \\
(26,14)\,5 & (26,16)\,7 & (26,18)\,5 & (26,20)\,7 & (26,22)\,5 & (26,24)\,7 \\
(26,26)\,5 & (28,2)\,281 & (28,4)\,31 & (28,6)\,11 & (28,8)\,17 & (28,10)\,13 \\
(28,12)\,19 & (28,14)\,31 & (28,16)\,11 & (28,18)\,631 & (28,20)\,109 & (28,22)\,13 \\
(28,24)\,17 & (28,26)\,11 & (30,2)\,5 & (30,4)\,17 & (30,6)\,5 & (30,8)\,11 \\
(30,10)\,5 & (30,12)\,17 & (30,14)\,5 & (30,16)\,373 & (30,18)\,5 & (30,20)\,17 \\
(30,22)\,5 & (30,24)\,83 \\
\end{longtable}
\end{document}
